\documentclass[11pt,letterpaper]{article}
\usepackage[letterpaper,margin=0.65in]{geometry}
\usepackage[utf8]{inputenc}
\usepackage[T1]{fontenc}
\usepackage[default]{lato}
\usepackage{enumitem}
\usepackage{titlesec}
\usepackage[hidelinks]{hyperref}

\hypersetup{
    pdftitle={Zhiqi Chen -- Curriculum Vitae},
    pdfauthor={Zhiqi Chen}
}
\urlstyle{same}
\pagestyle{plain}
\setlength{\parindent}{0pt}
\setlength{\parskip}{2pt}
\setlist[itemize]{leftmargin=1.2em,topsep=3pt,itemsep=2pt,parsep=0pt,partopsep=0pt}
\titleformat{\section}{\large\bfseries}{}{0pt}{}[\titlerule]
\titlespacing*{\section}{0pt}{10pt}{5pt}

\newcommand{\headerrow}[2]{%
    \begin{tabular*}{\linewidth}{@{}l@{\extracolsep{\fill}}r@{}}
        #1 & #2 \\
    \end{tabular*}\par
}

\begin{document}

\begin{center}
    {\LARGE\textbf{Zhiqi Chen}}\\[3pt]
    Department of Automation, Tsinghua University, Beijing, China\\[3pt]
    {\small
    \href{mailto:zq-chen23@mails.tsinghua.edu.cn}{zq-chen23@mails.tsinghua.edu.cn}
    \quad\textbar\quad
    \href{mailto:zqchen05@gmail.com}{zqchen05@gmail.com}
    \quad\textbar\quad +86 18367487119\\[2pt]
    \href{https://zhiqichen05.github.io/}{zhiqichen05.github.io}
    \quad\textbar\quad
    \href{https://github.com/zhiqichen05}{GitHub}
    \quad\textbar\quad
    \href{https://scholar.google.com/citations?user=Nn2pmyQAAAAJ}{Google Scholar}}
\end{center}

\section*{Education}

\headerrow{\textbf{Tsinghua University}}{Beijing, China}
\headerrow{B.Eng. in Automation, expected 2027}{Aug 2023 -- Present}
\begin{itemize}
    \item \textbf{GPA: 4.0/4.0; rank: 1/173} (as of July 2026).
    \item Selected coursework (all A+): Linear Algebra, Discrete Mathematics, Random Mathematics and Statistics, Operations Research, Data Structure, and Computer Languages and Programming.
    \item Member of the THU Tong Class (AGI Class), an interdisciplinary AI training program, since February 2024.
\end{itemize}

\section*{Research Interests}

I aim to understand how models \textbf{acquire new capabilities} and \textbf{use the ones they already have}, and to develop \textbf{training-data strategies and reinforcement learning methods} that expand reasoning and generalization in large language models. Interactive learning for agents, including robotic agents, is a complementary way to probe and extend these capabilities.

\section*{Publications}

{\small $^{*}$Equal contribution.}

\begin{enumerate}[leftmargin=1.5em,topsep=4pt,itemsep=7pt,parsep=0pt]
    \item
    \textbf{Does Reinforcement Learning Really Incentivize Reasoning Capacity in LLMs Beyond the Base Model?}\\
    Yang Yue$^{*}$, \textbf{Zhiqi Chen}$^{*}$, Rui Lu, Andrew Zhao, Zhaokai Wang, Yang Yue, Shiji Song, Gao Huang.\\
    \emph{NeurIPS 2025} (Oral), \textbf{Best Paper Runner-Up Award}.\\
    \emph{ICML 2025 AI4Math Workshop} (Oral), \textbf{Best Paper Award}.\\
    \href{https://arxiv.org/abs/2504.13837}{[Paper]} \href{https://github.com/LeapLabTHU/limit-of-RLVR}{[Code]}

    \item
    \textbf{Learning Push-Grasp Synergy for Occluded Objects in Cluttered Environments}\\
    Ziang Li, Haorui Wu, \textbf{Zhiqi Chen}, Haozhe Zhang, Yuzhe Huang, Changshui Zhang.\\
    \emph{ICRA 2026}.
    \href{https://doi.org/10.1109/ICRA57385.2026.11695822}{[Paper]}
\end{enumerate}

\section*{Research Reports}

\begin{enumerate}[leftmargin=1.5em,topsep=4pt,itemsep=7pt,parsep=0pt]
    \item
    \textbf{Code as Primitives: Harnessing Frontier Models to Synthesize Low-Level Robotic Skills}\\
    \textbf{Zhiqi Chen}$^{*}$, Djengo Cyun-Jyun Fang$^{*}$, Guanning Zeng, Xiuyu Li, Haiwen Feng, Andrea Zanette.\\
    \emph{Research blog, 2026}. \textbf{Project lead}.
    \href{https://zhiqichen05.github.io/Code-as-Primitives/}{[Blog]}
\end{enumerate}

\section*{Research Experience}

\headerrow{\textbf{Independent Research Project}}{Oct 2026 -- Present}
\textit{Project Lead} \hfill Advisor: Andrea Zanette\par
\textbf{Data Selection for Supervised Fine-Tuning}
\begin{itemize}
    \item Designing evaluation protocols and controlled baselines to study how data selection for supervised fine-tuning shapes reasoning models.
\end{itemize}

\headerrow{\textbf{Impossible Inc.}}{Jul 2026 -- Present}
\textit{Research Intern} \hfill Advisors: Andrea Zanette and Haiwen Feng\par
\textbf{Code as Primitives: Agentic Learning for Robotic Control}\par
\textit{Project Lead}
\begin{itemize}
    \item Led problem formulation, method design, implementation, and evaluation of a framework in which frontier-model agents synthesize low-level robot controllers from simulation feedback, without demonstrations.
    \item Proposed separating controller-program revision from numerical parameter search, so that each structural revision is based on focused, comparable evidence.
    \item Frozen controllers reached over 90\% held-out success across 75 DexVerse tasks and 99.2\% (119/120) across four Adroit tasks, with no model calls during execution.
\end{itemize}

\headerrow{\textbf{LeapLab, Tsinghua University}}{Feb 2025 -- Present}
\textit{Research Intern} \hfill Advisor: Prof. Gao Huang\par
\textbf{Limits of Reinforcement Learning with Verifiable Rewards}\par
\textit{Co-First Author}
\begin{itemize}
    \item Ran all main pass@$k$ comparisons between base and RL-trained models in the paper, spanning eleven benchmarks in math, code, and visual reasoning with up to 1,024 samples per problem.
    \item Conducted most of the paper's analyses, including perplexity analysis of reasoning traces, manual validity checks of chains of thought, and re-evaluation on filtered problems.
    \item These experiments provided the paper's core evidence that RLVR can improve single-sample accuracy while reducing coverage of solvable problems at large sampling budgets.
\end{itemize}

\headerrow{\textbf{Physical AI Research, Tsinghua University}}{Jul 2025 -- Sep 2025}
\textit{Research Intern} \hfill Advisor: Prof. Changshui Zhang\par
\textbf{Reinforcement Learning for Push--Grasp Manipulation}
\begin{itemize}
    \item Worked on reinforcement learning for coordinated pushing and grasping of occluded objects in cluttered environments.
    \item Deployed algorithms on physical robots for experimental validation, contributing to the work published at ICRA 2026.
\end{itemize}

\section*{Teaching Experience}

\headerrow{\textbf{Tsinghua University}}{Academic Year 2024 -- 2025}
\textit{Teaching Assistant, Calculus A}
\begin{itemize}
    \item Served for a full academic year as a second-year undergraduate, a role usually held by graduate students.
\end{itemize}

\section*{Honors and Awards}

\begin{itemize}
    \item \textbf{NeurIPS 2025 Best Paper Runner-Up Award}, as co-first author.
    \item \textbf{ICML 2025 AI4Math Workshop Best Paper Award}, as co-first author.
    \item \textbf{SenseTime Scholarship}, 2026; awarded to 30 undergraduates nationwide annually.
    \item \textbf{National Scholarship of China}, 2025 and 2024.
    \item \textbf{Person of the Year}, Department of Automation, Tsinghua University, 2025.
    \item \textbf{Huang-Yicong AGI Scholarship}, Tsinghua University, 2025.
    \item \textbf{Zheng-Weimin Scholarship, First Prize}, Tsinghua University, 2025.
\end{itemize}

\section*{Technical Skills}

\textbf{Machine learning:} SFT and RL training of large language models; large-scale sampling-based evaluation across model families; reimplementation of data-selection methods.\\
\textbf{Frameworks:} verl, vLLM, PyTorch, Hugging Face Transformers, Ray.\\
\textbf{Programming:} Python, C, C++. \quad \textbf{Scientific writing:} \LaTeX.

% \section*{Leadership and Activities}

% \textbf{Men's Football Team, Department of Automation, Tsinghua University}\\
% Captain, 2025--2026 season; player, 2023--2025 and 2026--2027 seasons, competing in the Tsinghua University Football League.

\end{document}
